\documentclass[11pt]{article}
\usepackage[margin=1in]{geometry}
\usepackage{amsmath,amssymb,amsthm}
\usepackage[hidelinks]{hyperref}
\newtheorem{theorem}{Theorem}
\newtheorem{lemma}{Lemma}
\newtheorem{corollary}{Corollary}
\newcommand{\cp}{\chi^{(2,1)}}
\newcommand{\cl}{\chi_{\ell}^{(2,1)}}
\title{A five-list obstruction on an eleven-vertex tree}
\author{Alec Kriebel\\
\small Independent researcher\\
\small \href{https://orcid.org/0009-0001-9320-500X}{ORCID: 0009-0001-9320-500X}}
\date{17 October 2026 \quad Version 1.0}
\begin{document}
\maketitle
\begin{abstract}
We exhibit an eleven-vertex tree \(T\) with ordinary \(L(2,1)\) initial-segment
palette size \(5\) and arbitrary-list parameter \(6\).
An explicit assignment of five-element positive-integer lists is infeasible
by a short forcing argument. An ordinary five-palette labeling and the known
\(\Delta+3\) tree list bound give the exact values.
This contradicts the universal equality stated in Kohl's 2006
Conjecture~4 and dissertation Conjecture~4.4.
Independent pendant-path extensions give infinitely many examples of
unbounded diameter. Historical priority, minimal order and present openness
are not claimed.
\end{abstract}

\section{Conventions and the historical target}
An \(L(2,1)\)-labeling of a finite simple graph assigns positive integers to
vertices so that adjacent labels differ by at least \(2\), and labels at
distance \(2\) differ. Write \(\cp(G)\) for the least \(k\) allowing such a
labeling from the common initial segment \([k]=\{1,\ldots,k\}\).
Write \(\cl(G)\) for the least \(k\) such that \emph{every} assignment of
\(k\)-element subsets \(L(v)\) of the positive integers admits a labeling
\(f(v)\in L(v)\). Lists need not be intervals.
Ordinary zero-based span \(\lambda(G)\) satisfies
\(\cp(G)=\lambda(G)+1\).

Kohl's Oberwolfach contribution \cite[Conjecture~4, p.~416]{kohl-owr}
and dissertation \cite[Conjecture~4.4, p.~124]{kohl-thesis}
state \(\cl(T)=\cp(T)\) for every tree \(T\).
The example below is a full negative answer to that stated universal equality.
The identification is historical; it does not assert that the conjecture
remained open immediately before this work.

\section{The obstruction}
Let \(T\) have the path \(a-b-c-d-e\). Attach leaves \(a_1,a_2\) to \(a\),
\(b_1\) to \(b\), \(d_1\) to \(d\), and \(e_1,e_2\) to \(e\).
These are all vertices and edges, so \(T\) has eleven vertices, ten edges,
maximum degree \(3\), radius \(3\), and diameter \(6\).
Set
\[
 A=\{1,2,3,4,5\},\qquad B=\{1,3,4,5,6\}.
\]
Assign \(B\) to \(a_1,a_2,e_1,e_2\) and \(A\) to every other vertex.

\begin{theorem}\label{thm:main}
The displayed assignment has no \(L(2,1)\)-labeling, while
\(\cp(T)=5\), or equivalently \(\lambda(T)=4\), and \(\cl(T)=6\).
\end{theorem}
\begin{proof}
In a putative list labeling, the three neighbors of \(b\) all have list \(A\)
and are pairwise at distance \(2\).
If \(f(b)=2,3,4\), respectively, at most two elements of \(A\)
differ from \(f(b)\) by at least \(2\). Thus \(f(b)\in\{1,5\}\).
The same argument applies to \(d\).
Since \(b,d\) are at distance \(2\), they receive different endpoint labels.
Reversing the spine and its attached leaves preserves the assignment,
so assume \(f(b)=1\) and \(f(d)=5\). Then \(f(c)=3\).

The edge \(ab\) and the distance-two pair \(ac\) force \(f(a)\in\{4,5\}\).
Each of \(a_1,a_2\) has list \(B\), must differ by at least \(2\) from \(f(a)\),
and must avoid \(1=f(b)\).
If \(f(a)=4\), only \(6\) remains; if \(f(a)=5\), only \(3\) remains.
In either case the two leaves receive the same label, although they are
at distance \(2\). This contradiction proves infeasibility.

An ordinary labeling from \([5]\) is
\[
\begin{array}{c|rrrrrrrrrrr}
 v & a&b&c&d&e&a_1&a_2&b_1&d_1&e_1&e_2\\ \hline
 f(v)&1&5&3&1&5&3&4&2&4&2&3
\end{array}
\]
as checked directly on edges and pairs of neighbors.
A degree-three vertex cannot be labeled from \([4]\): each possible center
label leaves at most two admissible neighbor labels, whereas its three
neighbors must be distinct. Hence \(\cp(T)=5\).
The bad five-lists imply \(\cl(T)\geq6\); Lemma~\ref{lem:bound} gives
the reverse inequality.
\end{proof}

\begin{lemma}[Known tree bound]\label{lem:bound}
For every nonempty finite tree \(S\),
\(\cl(S)\leq\Delta(S)+3\).
\end{lemma}
\begin{proof}
Root the tree and label in breadth-first order.
For a nonroot vertex \(v\), its parent is its only earlier neighbor,
excluding at most three integers from \(L(v)\).
Earlier vertices at distance \(2\) are its grandparent, if present,
and earlier siblings. There are at most \(\Delta(S)-1\) of these:
if the parent is the root there is no grandparent, and otherwise
the parent has at most \(\Delta(S)-1\) children.
They exclude at most one label each.
At most \(\Delta(S)+2\) labels are therefore forbidden, leaving a choice
from any \((\Delta(S)+3)\)-element list. The root has no earlier constraints.
This also covers a singleton.
\end{proof}
For positive integers \(d\geq s\), the more general tree bound
\(\chi_{\ell}^{(d,s)}(S)\leq s\Delta(S)+2d-1\)
is attributed to Voigt in \cite[p.~416 and reference~8 on p.~417]{kohl-owr},
citing Kohl, Schreyer, Tuza and Voigt \cite{kstv}.
The proof above is included for self-contained verification, without a
novelty claim for the bound.

\section{Infinitely many countertrees}
\begin{corollary}
Attach one path, of any independently chosen nonnegative finite length,
outward from each of the six leaves of \(T\).
Every resulting tree \(S\) has \(\cp(S)=5\) and \(\cl(S)=6\).
The family has unbounded diameter.
\end{corollary}
\begin{proof}
The old tree is an induced subtree, its old distances are unchanged, and
the maximum degree remains \(3\).
Extend the displayed ordinary labeling successively along each new path.
The parent excludes at most three labels of \([5]\), and the grandparent
at most one, so a choice always remains.
New vertices on different added paths have distance at least \(4\);
no further constraints arise. Thus \(\cp(S)\leq5\), with the reverse inequality
by restriction to \(T\).
Extend the infeasible lists on \(T\) by arbitrary five-element lists on
new vertices. Any labeling would restrict to a forbidden labeling of \(T\),
so \(\cl(S)\geq6\), while the tree bound gives \(\cl(S)\leq6\).
\end{proof}

\section{Verification, attribution and limitations}
The proof is elementary and independent of computational search.
The accompanying verification archive contains the exact graph, lists and
ordinary witness, TeX source, a standard-library Python enumerator,
expected results and byte manifests. Its exhaustive finite counts are
\(0,8,6336\) for common palettes \([4],[5],[6]\), respectively, and \(0\)
for the displayed five-lists. These computations check fixed instances;
the universal six-list bound and the entire infinite family follow from
the proofs, not from sampling.
A separate construction-input reconstruction used different constraint
propagation and tree-state mechanisms, froze its results before comparison,
and did not read the candidate proof. The construction was supplied to it,
so this is not an independent discovery claim.
Supplementary recurrence details include explicit local-state admissibility;
the main obstruction needs no recurrence.

General-graph separations were known earlier.
Fiala and \v{S}krekovski's 2001 preprint \cite[Fig.~4, p.~14]{fs}
displays a ten-vertex cyclic graph with an ordinary five-palette labeling
and infeasible five-lists.
That figure does not itself contradict tree equality.
A bounded primary-source and citation audit did not verify an earlier
negative tree resolution. Relevant inaccessible or unlocated source bodies
remain, as recorded in the supporting priority audit.
Accordingly this note makes no claim of historical firstness, minimality,
a new obstruction mechanism, or present openness.
The inspected 2001 preprint is cited as such; its content is not silently
identified with the differently titled 2005 journal version.

AI tools were used extensively in the research, proof development,
literature review, computation, verification and writing.
The internal independent and adversarial checks were AI-assisted.
This preprint is unrefereed: no external human peer review or formal
proof-assistant certification is claimed.

\begin{thebibliography}{9}
\bibitem{kohl-owr}
A. Kohl, Some notes on \(L(d,s)\)-list labellings of trees and cacti,
in \emph{Algorithmic Graph Theory}, Oberwolfach Reports \textbf{3} (2006), report 7,
414--417. Conjecture~4 on p.~416.
\href{https://doi.org/10.4171/OWR/2006/07}{doi:10.4171/OWR/2006/07}.

\bibitem{kohl-thesis}
A. Kohl, \emph{Knotenf\"arbungen mit Abstandsbedingungen},
doctoral dissertation, TU Bergakademie Freiberg, 2006,
\S4.4, especially p.~124. Institutional deposit dated 16 December 2009.
\href{https://nbn-resolving.org/urn:nbn:de:bsz:105-5454783}{urn:nbn:de:bsz:105-5454783}.

\bibitem{kstv}
A. Kohl, J. Schreyer, Z. Tuza and M. Voigt,
List version of \(L(d,s)\)-labelings,
\emph{Theoretical Computer Science} \textbf{349}(1) (2005), 92--98.
\href{https://doi.org/10.1016/j.tcs.2005.09.032}{doi:10.1016/j.tcs.2005.09.032}.

\bibitem{fs}
J. Fiala and R. \v{S}krekovski,
\emph{List Distance Labelings of Graphs},
KAM-DIMATIA Series 2001-530; ITI Series 2001-28, 2001.
\href{https://kam.mff.cuni.cz/kamserie/clanky/2001/s530.pdf}{Institutional preprint},
Fig.~4, p.~14.
\end{thebibliography}
\end{document}
